\item Un sistema que consiste en $n$ moles de un gas ideal con calor específico molar a presión constante $C_P$ experimenta dos procesos reversibles. Comienza con presión $P_i$ y volumen $V_i$, se expande isotérmicamente y luego se contrae adiabáticamente para alcanzar un estado final con presión $P_i$ y volumen $3V_i$. (a) Encuentre su cambio de entropía en el proceso isotérmico. (La entropía no cambia en el proceso adiabático.) (b) \textbf{¿Qué pasaría si?} Explique por qué la respuesta al inciso (a) debe ser la misma que la respuesta al problema 46. (No es necesario resolver el problema 46 para contestar esta pregunta.)
